\section{Limitations}
\label{sec:limits}

We state limitations plainly and do not minimize them.

\paragraph{Training budget and design.} Corpora are byte-matched at a small
$32$\,MiB budget, bounded by the available cleaned Azure training text; absolute
token counts are budget-dependent. The controlled matrix uses a single tokenizer
algorithm (byte-level BPE). Sources are technical documentation, so results may
not extend to chat, source-code repositories, prose, or all enterprise text.

\paragraph{The Azure-versus-general-English confound (most important).} Comparing
Azure documentation with general English changes more than subject domain; it also
changes technical writing style, markup distribution, identifier density, and
document structure. The transfer and prose-only analyses narrow this concern but
do \emph{not} fully isolate Azure-specific vocabulary from general
technical-documentation specialization. A matched general-technical documentation
training control is the priority follow-up (Section~\ref{sec:future}).

\paragraph{Scope of measurement.} We train no language model and evaluate no
downstream task; there is no GPU or serving benchmark. CPU runtime was measured on
a thermally variable laptop, so small runtime differences are not overinterpreted,
and CPU tokenizer throughput is not model-serving throughput.

\paragraph{Data and comparison caveats.} Results are tied to pinned public-source
snapshots that change over time; raw corpora are not redistributed, and
reproduction downloads the pinned sources. Existing production tokenizers (Track
B) are observational, not controlled. Any context-window or token-billing figure
elsewhere is a derived estimate, not a measured GPU or billing outcome. We claim
no universal domain-tokenizer superiority.

\paragraph{Affiliation.} If the author is affiliated with the vendor of the studied
documentation, that relationship is disclosed where applicable; this study uses
only public documentation, is not an official result of that vendor, and implies
no endorsement.
